% Imporditud: tools/import_eksperimendid.py
% Allikas: Eesti füüsikaolümpiaadi eksperimendi ülesannete
%   kogu (Rael Kalda, 2023), ülesanne Ü139, lahendus L139.
\setAuthor{EFO žürii}
\setRound{lõppvoor}
\setYear{2010}
\setNumber{G E2}
\setDifficulty{4}
\setTopic{dünaamika}
\setVahendid{Tuntud massiga õhuke plaat (m = 20 g), statiiv, joonlaud, ventilaator, niit}
\setPunktid{12}
\setAllikas{Eksperimendi kogu Ü139, lahendus lk 104}
\setLahendusseis{toores}

\prob{Ventilaator}
Leidke hinnang ventilaatorist väljuva õhujoa kiiruse jaoks. Õhu tihedus $\rho$ = 1,29 kg/m$^3$, raskuskiirendus g = 9,8 m/s$^2$.

\hint

\solu
Riputame lauakese niidi abil ühest otsast kinnitades statiivi külge ning suuname

ventilaatorist tuleva õhuvoo laua alumise osa vastu. Mõõdame õhuvoolust tingi-

tud laua alumise otsa kõrvalekalde a. Ajaühikus laua alumisele otsale langeva õhu mass m/t = $\pi$$d_2$$\rho$v/4 (kus d on ventilaatori diameeter) ning see õhumass kannab impulssi mv/t = $\pi$$d_2$$\rho$$v_2$/4, mis antakse üle lauale ja kujutab endast efektiivset jõudu. Selle jõu moment tasakaalustab raskusjõu momendi M ga/2, st $\pi$Ld2$\rho$$v_2$/4 = M ga/2, millest v = 2M ga/$\pi\rho$L/d.

Tulemuseks saame v = 3,8 m/s.
\probend
